\section{Preserving wealth}

\begin{qfmethod}[Risks to human capital and how to manage them: mortality, disability, unemployment, wage and industry risk, correlation with financial assets, and health, with insurance, diversification and reserves]\label{meth:human-capital-risks}
\qffield{Statement} \emph{Types of risk to human capital} (Interpretation, extending \citep[pp.~13--16, 28--31, 75--76]{ibbotson-lifetime-financial-advice-2007}):
\begin{itemize}
\item \emph{mortality}: earnings stop for the family;
\item \emph{disability and ill health}: earnings stop or fall while costs rise;
\item \emph{unemployment and wage risk}: lost or lower pay, often in recessions when financial assets also fall;
\item \emph{industry and employer risk}: skills tied to a declining industry or a single employer, worst when the same employer's stock is in the portfolio;
\item \emph{obsolescence}: skills that lose value as technology changes.
\end{itemize}
Human capital's correlation with markets decides how much financial risk the household can carry: equity-like earnings call for more bonds in the financial portfolio; bond-like earnings allow more equities (pp.~14--16, 76).

\emph{Strategies}: life insurance for mortality (term cover sized to human capital not otherwise funded); disability income cover; health and long-term-care cover; an emergency reserve sized to income volatility; diversification away from employer stock and the client's own industry; investment in skills; and, for business owners, key-person cover and succession plans (Interpretation).

\qffield{Practice} \emph{Tunde Okafor} (original; \texttt{private-wealth-checks.py}, parts~6 and 10). Tunde, 42, earns mainly commission from a listed software company, and holds 40 percent of the family's liquid assets in its shares (volatility 45 percent, correlation 0.5 with a 16 percent-volatility market).
\begin{itemize}
\item The concentrated portfolio's volatility is 24.3 percent against 16 percent diversified; and its worst outcomes coincide with his commission income falling.
\item \emph{Analysis}: his human capital is equity-like and tied to the same company as 40 percent of the portfolio, so a downturn there hits income, employment and savings together.
\item \emph{Recommendations}: sell the shares down to below 10 percent of liquid assets over a year (a pre-arranged Rule~10b5-1 plan if he is an insider); hold a larger emergency reserve (six months) than a salaried employee would; keep the household's financial portfolio more bond-heavy than Adaeze's bond-like income alone would justify; buy disability cover that includes commission income; and term life cover on both spouses, since Adaeze's 2.17 million need (above) is matched by a smaller one for Tunde.
\end{itemize}

\qffield{Conditions} Volatilities and correlation are illustrative; insurance terms and taxation of benefits differ by policy and jurisdiction.

\qffield{Pitfalls} Holding employer stock because it is familiar. Insuring only the higher earner. Ignoring the correlation between job loss and market falls.

\qffield{See also} \hyperref[meth:lifetime-exposures]{lifetime exposures}; \hyperref[meth:complex-client-situations]{executives}.

\end{qfmethod}

\begin{qfmethod}[Preserving purchasing power: inflation-protection strategies for spending and investments, and how exchange rates affect allocation and plans, with ways to reduce the risk]\label{meth:purchasing-power-fx}
\qffield{Statement} \emph{Inflation.} Spending needs are real, but many assets and incomes are nominal. A level nominal pension or annuity loses purchasing power steadily; taxes on nominal returns make it worse (\hyperref[meth:private-client-tax-inflation]{taxes and inflation}). Strategies: inflation-linked bonds such as TIPS, whose principal is indexed to CPI-U \citep[web page]{treasurydirect-tips}; equities and real assets, which tend to protect over long horizons but not reliably over short ones; inflation-indexed or escalating annuities and state pensions; holding inflation-linked assets in tax-advantaged accounts; and planning spending in real terms with a margin (Interpretation, with \hyperref[def:tips-breakevens]{TIPS}).

\emph{Exchange rates.} A client who spends in one currency and holds assets in others bears currency risk on those assets. Hedging foreign bonds removes most of their risk, since currency returns are almost uncorrelated with bond returns; for global equities some currencies (in 1975--2005 the US dollar, euro and Swiss franc) moved against equity markets and reduced risk, so a full hedge is not always risk-minimising \citep[NBER WP 13088, pp.~1--3]{campbell-viceira-currency-hedging-2007}. Practical steps: match the currency of assets to the currency of planned spending, especially for near-term goals; hedge foreign bonds; decide a hedge ratio for foreign equities; and for clients with multi-currency lives (expatriates, retirees abroad), hold separate currency buckets for each spending stream (Interpretation; \hyperref[meth:currency-hedging-programme]{currency programme}).

\qffield{Practice} (original; \texttt{private-wealth-checks.py}, part~10).
\begin{itemize}
\item \emph{Inflation}: a level nominal pension of 50,000 is worth 27,700 in today's money after 20 years of 3 percent inflation, a 45 percent loss. For the Okafors' essential retirement spending above Social Security, a TIPS ladder bought gradually in their 401(k)s matches real spending and avoids tax on the inflation uplift.
\item \emph{Exchange rates}: Margaret Hartley lives on sterling with 30 percent of her portfolio in US dollar assets. A 10 percent rise in sterling cuts the portfolio's sterling value by 2.7 percent. Recommendation: hedge her dollar bonds fully and half her US equities, keeping the rest unhedged for diversification, and hold three years of spending in sterling assets.
\end{itemize}

\qffield{Conditions} Inflation and currency figures are illustrative; Campbell and co-authors' results are for 1975--2005 and a one-quarter horizon.

\qffield{Pitfalls} Treating nominal bonds as safe for real goals. Hedging all foreign equities by reflex. Ignoring the currency of future spending for internationally mobile clients.

\qffield{See also} \hyperref[meth:lifetime-exposures]{longevity}; \hyperref[meth:residency-citizenship]{residency}.

\end{qfmethod}
